\section{ÉLECTROSTATIQUE}

% ============================================================
\subsection{Charge électrique}
\begin{formules}
  \ligne{Charge électrique}
    {$q,\ e,\ Q$}
    {C}
    {Charge élémentaire : \newline $e = 1.602\,176\,634\times10^{-19}$\,C}
  \ligne{Permittivité}
    {$\varepsilon = \varepsilon_r\varepsilon_0$}
    {F/m}
    {$\varepsilon_0$ : permittivité du vide \newline
     $\varepsilon_0 = 8.854187\times10^{-12}\ \frac{\mathrm{A\cdot s}}{\mathrm{V\cdot m}}$ \newline
     $\varepsilon_r$ : permittivité relative (constante diélectrique), \newline
     $[\varepsilon_r] = 1$ ; vide : $\varepsilon_r = 1$}
  \ligne{Matériau diélectrique (linéaire, isotrope, homogène)}
    {$\varepsilon_r = 1 + \chi_e$}
    {[--]}
    {$\chi_e$ : susceptibilité électrique \newline [--]}
  \ligne{Loi de Coulomb}
    {$\begin{gathered}
      \text{Scalaire : } F_{12} = \frac{1}{4\pi\varepsilon}\,\frac{q_1q_2}{r_{12}^2} \\
      \text{Vectorielle :} \\
      \vec F_{12} = \frac{1}{4\pi\varepsilon}\,\frac{q_1q_2}{r_{12}^2}\,\frac{\vec r_{12}}{r_{12}} \\
      \vec F_{12} = \frac{1}{4\pi\varepsilon}\,\frac{q_1q_2}{r_{12}^3}\,\vec r_{12} \\
      \textcolor{red}{\scriptstyle F_{12}\ :\ 1\,=\,\text{émet},\ 2\,=\,\text{subit}}
     \end{gathered}$}
    {N}
    {\schema{
      \node at (0.4,0.62) {même signe};
      \fill (0,0) circle (0.9pt) node[below] {$q_1$};
      \draw[->] (0,0) -- (0.48,0.24) node[midway, above left] {$\vec r_{12}$};
      \fill (0.5,0.25) circle (0.9pt) node[below right] {$q_2$};
      \draw[->, red, thick] (0.5,0.25) -- (0.85,0.43) node[right, red] {$\vec F_{12}$};
      \node at (1.95,0.62) {signes opposés};
      \fill (1.55,0) circle (0.9pt) node[below] {$q_1$};
      \draw[->] (1.55,0) -- (2.03,0.24);
      \fill (2.05,0.25) circle (0.9pt) node[right] {$q_2$};
      \draw[->, red, thick] (2.05,0.25) -- (1.8,0.125);
      \node[red] at (1.95,-0.02) {$\vec F_{12}$};
     }
     $q_1, q_2$ : charges ; $\vec F_{12}$ : force de Coulomb exercée par $q_1$ sur $q_2$ ;
     $\vec r_{12}/r_{12}$ : vecteur unitaire. Sens donné par le signe de $q_1q_2$ :
     $q_1q_2 > 0$ répulsion, $q_1q_2 < 0$ attraction}
  \ligne{Densité linéique}
    {$\lambda = \dfrac{q}{\ell}$}
    {C/m}
    {$\ell$ : longueur}
  \ligne{Densité surfacique}
    {$\sigma = \dfrac{q}{S}$}
    {C/m$^2$}
    {$S$ : surface}
  \ligne{Densité volumique}
    {$\rho = \dfrac{q}{V}$}
    {C/m$^3$}
    {$V$ : volume}
\end{formules}

% ============================================================
\subsection{Champ électrique}
\begin{formules}
  \ligne{Champ électrique}
    {$\vec F = q\vec E$}
    {N}
    {$\vec E$ : champ électrique ; $q$ : charge plongée dans $\vec E$ ; $\vec F$ : force exercée sur $q$}
  \soustitre{Champs électriques particuliers}
  \ligne{Intérieur d'un conducteur}
    {$E = 0$}
    {V/m}
    {}
  \ligne{Charge ponctuelle ou conducteur sphérique}
    {$\begin{gathered}
      \text{Scalaire : } E = \frac{1}{4\pi\varepsilon}\,\frac{q}{r^2} \\
      \text{Vectorielle :} \\
      \vec E = \frac{1}{4\pi\varepsilon}\,\frac{q}{r^2}\,\frac{\vec r}{r} \\
      \vec E = \frac{1}{4\pi\varepsilon}\,\frac{q}{r^3}\,\vec r \\
      \text{Condition : } r > R
     \end{gathered}$}
    {V/m}
    {\schema{
      \fill (0,0) circle (1pt) node[below] {$q > 0$};
      \draw[->] (0,0) -- (0.65,0.23) node[midway, above] {$\vec r$};
      \fill (0.67,0.24) circle (0.8pt) node[below right] {P};
      \draw[->, champ, thick] (0.67,0.24) -- (1.15,0.41) node[right] {$\vec E$};
      \node[anchor=west] at (0.75,-0.02) {sens opposé si $q < 0$};
     }
     $\vec E$ : champ créé par la charge $q$ au point P ; $R$ : rayon de la sphère}
  \ligne{Ensemble de charges ponctuelles}
    {$\vec E = \sum \vec E_i$}
    {V/m}
    {$\vec E$ : champ résultant de la somme des champs $\vec E_i$}
  \ligne{Conducteur cylindrique}
    {$\begin{gathered}
      E = \frac{1}{2\pi\varepsilon}\,\frac{Q}{rh} \\
      \text{Condition : } r > R
     \end{gathered}$}
    {V/m}
    {\schema{
      \draw[<->] (0,0.62) -- (1.6,0.62) node[midway, above] {$h$};
      \draw (0,0.42) -- (0,0.66) (1.6,0.42) -- (1.6,0.66);
      \draw (0.05,0.1) -- (1.55,0.1) (0.05,0.38) -- (1.55,0.38);
      \draw (0.05,0.24) ellipse (0.06 and 0.14);
      \draw (1.55,0.24) ellipse (0.06 and 0.14);
      \node at (0.8,0.24) {$Q > 0$};
      \draw (0.8,0.1) -- (0.8,-0.25);
      \fill (0.8,-0.25) circle (0.8pt) node[left] {P};
      \draw[->, champ, thick] (0.8,-0.25) -- (0.8,-0.75) node[right] {$\vec E$};
     }
     $\vec E$ : champ créé par le cylindre au point P ; $h$ : hauteur du cylindre ;
     $R$ : rayon du cylindre. Effets de bord négligés (cylindre infini). Sens opposé si $Q < 0$}
  \ligne{Conducteur plan}
    {$E = \dfrac{1}{2\varepsilon}\,\dfrac{Q}{S}$}
    {V/m}
    {\schema{
      \draw[fill=gray!8] (0,0) -- (0.3,0.22) -- (0.3,1.0) -- (0,0.78) -- cycle;
      \node at (0.1,0.7) {$S$};
      \draw[->, champ, thick] (0.2,0.5) -- (1.0,0.5) node[right] {$\vec E$};
      \node[anchor=west] at (-0.05,-0.12) {$Q > 0$};
      \node[anchor=west] at (0.75,-0.12) {sens opposé si $Q < 0$};
     }
     $\vec E$ : champ créé par le plan en tous points ; $S$ : surface du plan.
     Effets de bord négligés (plan infini)}
  \ligne{Entre deux conducteurs plans}
    {$E = \dfrac{1}{\varepsilon}\,\dfrac{Q}{S}$}
    {V/m}
    {\schema{
      \fill[champ!12] (0.66,0) rectangle (1.44,0.9);
      \draw[fill=gray!25] (0.58,0) rectangle (0.66,0.9);
      \draw[fill=gray!25] (1.44,0) rectangle (1.52,0.9);
      \foreach \y in {0.15,0.35,0.55,0.75} \draw[->, champ] (0.7,\y) -- (1.4,\y);
      \node[champ] at (1.05,1.02) {$\vec E$};
      \node at (0.3,0.75) {$E\!=\!0$};
      \node at (1.82,0.75) {$E\!=\!0$};
      \node[below] at (0.62,0) {$+Q$};
      \node[below] at (1.48,0) {$-Q$};
     }
     $S$ : surface de recouvrement des deux plans. Effets de bord négligés (plans infinis)}
  \ligne{Champ de déplacement}
    {$\begin{gathered}
      \vec D = \varepsilon\vec E \\
      D = \varepsilon E = \frac{Q}{S}
     \end{gathered}$}
    {C/m$^2$}
    {$\vec D$ : champ de déplacement ou induction électrique}
  \ligne{Flux élémentaire}
    {$\vec E\cdot d\vec A = E\,dS\cos\theta$}
    {}
    {$\theta$ : angle entre $\vec E$ et la normale $d\vec A$}
  \ligne{Théorème de Gauss}
    {$\begin{gathered}
      \oiint_S \vec E\cdot d\vec S = \frac{1}{\varepsilon}\sum_i q_i = \frac{Q_{int}}{\varepsilon_0} \\
      = \frac{1}{\varepsilon}\iiint_V \rho\,dV \\
      \text{Loi locale : } \vec\nabla\cdot\vec E = \frac{\rho}{\varepsilon} \\
      \vec E = -\vec\nabla V
     \end{gathered}$}
    {}
    {\schema{
      \draw (1.1,0.1) circle (0.42);
      \node at (1.25,0.3) {$V,\ \rho$};
      \draw[fill=white, rotate around={45:(0.8,0.4)}] (0.74,0.34) rectangle (0.86,0.46);
      \draw[->, thick] (0.8,0.4) -- (0.42,0.62) node[left] {$d\vec S$};
      \draw[->, champ, thick] (0.8,0.4) -- (0.55,0.05) node[left] {$\vec E$};
     }
     $S$ : surface fermée ; $d\vec S$ : élément de surface ; $q_i$ : charges contenues dans $S$ ;
     $V$ : volume contenu dans $S$ ; $\rho$ : densité de charges dans $V$ ; $\vec\nabla$ : opérateur nabla}
  \ligne{Coordonnées cylindriques}
    {$\begin{gathered}
      x = r\cos(\theta) \\
      y = r\sin(\theta) \qquad dS = r\,dz\,d\theta \\
      z = z
     \end{gathered}$}
    {}
    {}
  \ligne{Coordonnées sphériques}
    {$\begin{gathered}
      x = \rho\cos(\theta)\sin(\varphi) \\
      y = \rho\sin(\theta)\sin(\varphi) \\
      z = \rho\cos(\varphi) \\
      dS = \rho^2\sin(\varphi)\,d\varphi\,d\theta
     \end{gathered}$}
    {}
    {}
\end{formules}

% ============================================================
\subsection{Potentiel}
\begin{formules}
  \ligne{Loi générale}
    {$\begin{gathered}
      V = -\int_O^P \vec E\cdot d\vec r \\
      V_B - V_A = -\int_A^B \vec E\cdot d\vec r \\
      \lim_{r\to+\infty} V(r) = 0
     \end{gathered}$}
    {V}
    {$V$ : potentiel au point P ; O : référence des potentiels ; $d\vec r$ : élément de parcours de O à P}
  \ligne{Loi locale}
    {$\vec E = -\vec\nabla V$}
    {V/m}
    {}
  \ligne{Équipoten\-tielles}
    {Cercles (sphères) au même potentiel autour de $q$}
    {}
    {\schema{
      \draw[blue!60, dashed] (0,0) circle (0.28);
      \foreach \a in {0,45,...,315} \draw[->, champ] (\a:0.06) -- (\a:0.45);
      \fill (0,0) circle (1pt);
      \node[blue!70, anchor=west] at (0.55,0.15) {équipot.};
      \node[champ, anchor=west] at (0.55,-0.12) {ligne champ};
     }
     Lignes de champ $\perp$ équipotentielles}
  \ligne{Potentiel d'une charge ponctuelle}
    {$V = \dfrac{1}{4\pi\varepsilon}\,\dfrac{q}{r}$}
    {V}
    {$r$ : distance de $q$ au point P}
  \ligne{Tension ou différence de potentiel}
    {$\begin{gathered}
      U_{AB} = V_A - V_B = \int_A^B \vec E\cdot d\vec r \\
      \text{Plans : } U = \frac{q}{S}\,\frac{d}{\varepsilon_0}
     \end{gathered}$}
    {V}
    {$U_{AB}$ : tension entre les points A et B}
  \ligne{Travail}
    {$\begin{gathered}
      W_{AB} = qU_{AB} \\
      W = \vec F\cdot\vec d \quad W = Fx + Uq
     \end{gathered}$}
    {J}
    {$W_{AB}$ : travail à fournir à la charge $q$ pour la déplacer du point A au point B}
  \ligne{Énergie potentielle}
    {$E_p = QV$}
    {J}
    {}
\end{formules}

% ============================================================
\subsection{Capacité}
\begin{formules}
  \ligne{Capacité conducteur}
    {$C = \dfrac{Q}{V}$}
    {F}
    {$Q$ : charge du conducteur ; $V$ : potentiel du conducteur}
  \ligne{Capacité condensateur}
    {$C = \dfrac{Q}{U} \qquad Q = CU$}
    {F \newline C}
    {$Q$ : charge des armatures ($Q_+ = Q$, $Q_- = -Q$) ; $U$ : tension entre les armatures}
  \soustitre{Capacités de condensateurs particuliers}
  \ligne{Condensateur plan}
    {$C = \varepsilon\dfrac{S}{d}$}
    {F}
    {$S$ : surface de recouvrement des armatures ; $d$ : distance entre les armatures}
  \ligne{Condensateur sphérique}
    {$C = 4\pi\varepsilon\dfrac{1}{\frac{1}{r_1} - \frac{1}{r_2}}$}
    {F}
    {$r_1, r_2$ : rayons des sphères concentriques ($r_1 < r_2$)}
  \ligne{Sphère seule}
    {$C = 4\pi\varepsilon R$}
    {F}
    {$R$ : rayon de la sphère}
  \ligne{Condensateur cylindrique}
    {$C = 2\pi\varepsilon\dfrac{h}{\ln\frac{r_2}{r_1}}$}
    {F}
    {$r_1, r_2$ : rayons des cylindres concentriques ($r_1 < r_2$) ; $h$ : hauteur en vis-à-vis}
  \ligne{Ligne bifilaire}
    {$C = \dfrac{\pi\varepsilon L}{\ln\left(\frac{2d}{a} - 1\right)}$}
    {F}
    {$L$ : longueur de la ligne ; $d$ : distance entre les centres des fils ; $a$ : diamètre des fils}
  \soustitre{Combinaison de capacités}
  \ligne{En série}
    {$\dfrac{1}{C} = \sum\limits_i \dfrac{1}{C_i}$}
    {F}
    {}
  \ligne{En parallèle}
    {$C = \sum\limits_i C_i$}
    {F}
    {}
  \ligne{Énergie dans un condensateur}
    {$W = \frac{1}{2}CU^2$}
    {J}
    {}
  \ligne{Densité d'énergie électrique}
    {$w = \frac{1}{2}\varepsilon E^2 = \dfrac{W}{\text{Volume}}$}
    {J/m$^3$}
    {}
\end{formules}

% ============================================================
\subsection{Dipôle électrique}
\begin{formules}
  \ligne{Moment dipolaire}
    {$\vec p = q\vec d$}
    {C\,m}
    {\schema{
      \fill (0,0) circle (1pt) node[below] {$-q$};
      \draw[->, thick, blue!70!black] (0,0) -- (1.0,0.25) node[midway, above] {$\vec d$};
      \fill (1.02,0.255) circle (1pt) node[right] {$+q$};
     }
     $q$ : charge du dipôle ; $\vec d$ : de la charge négative à la positive}
  \ligne{Couple dans un champ $\vec E$}
    {$\vec M = \vec p\times\vec E$}
    {N\,m}
    {}
  \ligne{Énergie potentielle}
    {$W = \vec p\cdot\vec E$}
    {J}
    {Par convention, $W$ nulle quand $\vec d\perp\vec E$}
\end{formules}

% ============================================================
\subsection{Applications}
\begin{formules}
  \ligne{Barreau rectiligne infini uniformément chargé}
    {$\begin{gathered}
      E(r) = \frac{\lambda}{2\pi\varepsilon_0 r} \\
      \varepsilon = \varepsilon_0\varepsilon_r
     \end{gathered}$}
    {V/m}
    {Champ à une distance $r$ du barreau ; $\lambda$ : densité linéique}
  \ligne{Plaque infinie uniformément chargée}
    {$E = \dfrac{\sigma}{2\varepsilon_0}$}
    {V/m}
    {Ne dépend pas de la distance à la plaque !}
  \ligne{Condensateur plan}
    {$E = \dfrac{\sigma}{\varepsilon} = \dfrac{U}{d}$}
    {V/m}
    {Champ entre les armatures}
  \ligne{Polarisation}
    {$P = \chi_e\,\varepsilon_0\,E$}
    {C/m$^2$}
    {$\chi_e$ : susceptibilité électrique}
  \ligne{Rigidité diélectrique}
    {$E_{max}$}
    {V/m}
    {Champ maximal avant claquage du diélectrique}
  \ligne{Fil semi-fin}
    {$\begin{gathered}
      \vec E = \frac{1}{4\pi\varepsilon_0}\,\frac{q}{r^3}\,\vec r \\
      = \frac{1}{4\pi\varepsilon_0}\,\frac{q}{(x^2+y^2)^{3/2}}\begin{pmatrix}-x\\ y\end{pmatrix}
     \end{gathered}$}
    {}
    {}
  \ligne{Déplacement dans un champ}
    {$\begin{gathered}
      y(t) = \frac{1}{2}at^2 + v_0t + y_0 \\
      \vec F = q\vec E = m\vec a \ \Rightarrow\ a = \frac{qE}{m} \\
      v(t) = at + v_0
     \end{gathered}$}
    {}
    {}
\end{formules}
